\section{Factory Method}
\label{sec:pat-factory-method}

\problemstatement{Define an interface for creating an object, but let the
decision of \emph{which} concrete class to instantiate be made elsewhere,
so client code depends only on the abstract product, never on concrete
class names.}

\begin{patternbox}[When You Reach for It]
The trigger is \emph{``create an object without hard-coding its class,''}
usually to kill a \texttt{switch}/\texttt{if-else} on a type string that
\texttt{new}s different concrete classes. Whenever adding a new product
type forces you to edit existing creation code, a factory restores the
Open/Closed Principle.
\end{patternbox}

\subsection*{Understanding the pattern}

Client code should not contain \texttt{new ConcreteA} or \texttt{new
ConcreteB}. Instead it calls a factory function that takes a request
(often an enum or string) and returns a pointer to the abstract product.
The concrete decision lives in one place; the rest of the system speaks
only the interface. Adding a product means extending the factory, not
hunting down every \texttt{new} scattered across the codebase.

\begin{ideabox}[One Place Knows the Concrete Classes]
Centralise object creation behind a function returning the base type. The
caller says \emph{what} it wants (``a PDF exporter''); the factory decides
\emph{how} to build it. This single indirection is what lets you swap or
add implementations without touching callers.
\end{ideabox}

\subsection*{C++ implementation}

\begin{lstlisting}
#include <bits/stdc++.h>
using namespace std;

struct Shape {
    virtual void draw() const = 0;
    virtual ~Shape() = default;
};
struct Circle : Shape { void draw() const override { cout << "Circle\n"; } };
struct Square : Shape { void draw() const override { cout << "Square\n"; } };

enum class ShapeType { Circle, Square };

struct ShapeFactory {                          // the one place that knows
    static unique_ptr<Shape> create(ShapeType t) {
        switch (t) {
            case ShapeType::Circle: return make_unique<Circle>();
            case ShapeType::Square: return make_unique<Square>();
        }
        return nullptr;
    }
};

int main() {
    auto s1 = ShapeFactory::create(ShapeType::Circle);
    auto s2 = ShapeFactory::create(ShapeType::Square);
    s1->draw();                                // caller knows only "Shape"
    s2->draw();
    return 0;
}
\end{lstlisting}

\begin{complexitybox}[Trade-offs]
\textbf{Pros.} Decouples callers from concrete classes; new products are
added in one spot; supports the Open/Closed Principle. \textbf{Cons.} Adds
a layer of indirection and one more class/function to maintain; the
factory itself still centralises knowledge of every product (mitigated by
registration-based factories).
\end{complexitybox}

\begin{takeawaybox}
Factory Method replaces scattered \texttt{new ConcreteClass} calls with a
single creator returning the abstract type. It is the standard cure for a
type-tag \texttt{switch} and the most common way an interviewer expects
you to satisfy Open/Closed. Abstract Factory (next) generalises it to
whole \emph{families} of products.
\end{takeawaybox}
